\documentclass[../../../include/open-logic-section]{subfiles}

\begin{document}

\olfileid{sth}{ordinals}{iso}
\olsection{క్రమ-సమరూపతలు}

సుక్రమం \emph{ఎంత} దృఢమైన
భావమో వివరించడానికి,
సుక్రమాలను పోల్చే ఒక
పద్ధతిని పరిచయం చేస్తాం.

\begin{defn}
ఒక \emph{సుక్రమం} అంటే
$<$, $A$ను సుక్రమితం
చేసే $\tuple{A, <}$ జంట.
$\tuple{A, <}$,
$\tuple{B,
\lessdot}$ సుక్రమాలు
\emph{క్రమ-సమరూపమైనవి}
అంటే, $x < y$ అప్పుడే
మరియు అప్పుడే $f(x) \lessdot  f(y)$
అయ్యే ఒక
\tetoken{ద్విగుణ ప్రమేయం}{!!a{bijection}}
$f \colon A \to B$
ఉంటుంది. ఈ సందర్భంలో
$\ordeq{\tuple{A, <}}{\tuple{B, \lessdot}}$
అని రాసి, $f$ను
\emph{క్రమ-సమరూపత}
అంటాం.
\end{defn}
\noindent
ఇకపై సంక్షిప్తత కోసం
``క్రమ-సమరూపతలు'' బదులు
``సమరూపతలు'' అంటాం.
సహజంగా, సమరూపతలు
నిర్మాణాన్ని సంరక్షించే
\tetoken{ద్విగుణ ప్రమేయాలు}{!!{bijection}s}.
వాటి గురించి కొన్ని
సరళమైన విషయాలు ఇవి.

\begin{lem}\ollabel{isoscompose}
సమరూపతల సంయుక్తం
సమరూపతే; అంటే,
$f \colon A
\to B$, $g \colon B \to C$
సమరూపతలైతే,
$(g \circ f)
\colon A \to C$
సమరూపత అవుతుంది.
\end{lem}
\begin{prob}
	\olref[sth][ordinals][iso]{isoscompose}ను నిరూపించండి.
\end{prob}
\begin{proof}
అభ్యాసంగా వదిలాం.
\end{proof}
\begin{cor}\ollabel{ordisoisequiv}
	$\ordeq{X}{Y}$ ఒక
	సమానతా సంబంధం.
\end{cor}
\begin{prop}\ollabel{ordisounique}
$\tuple{A, <}$,
$\tuple{B, \lessdot}$
సుక్రమాలు సమరూపమైనవైతే,
వాటి మధ్య సమరూపత
ఏకైకం.
\end{prop}

\begin{proof}
$f$, $g$లు $A \to B$
సమరూపతలు అనుకుందాం.
\olref[wo]{propwoinduction}ను
ఉపయోగించి ఆగమనం
ద్వారా ఫలితాన్ని
నిరూపిస్తాం. $a\in A$ను
స్థిరపరచి, ఆగమన
పరికల్పనగా
$(\forall b < a)f(b) = g(b)$
అనుకుందాం. $x \in B$ను
స్థిరపరచండి.

$x \lessdot f(a)$ అయితే,
$f^{-1}(x) < a$; $f$,
$g$ సమరూపతలనే
విషయాన్ని ఉపయోగిస్తే
$g(f^{-1}(x)) \lessdot
g(a)$. కానీ
$f^{-1}(x) < a$
కాబట్టి మన పరికల్పన
ప్రకారం
$x =f(f^{-1}(x)) = g(f^{-1}(x))$.
అందువల్ల $x \lessdot g(a)$.
అదే విధంగా
$x \lessdot g(a)$ అయితే
$x \lessdot f(a)$.

సాధారణీకరిస్తే,
$(\forall x \in B)(x \lessdot f(a) \liff x \lessdot
g(a))$. అందువల్ల
\olref[sfr][rel][ord]{prop:extensionality-strictlinearorders}
ప్రకారం $f(a) = g(a)$.
కాబట్టి
\olref[wo]{propwoinduction}
ప్రకారం
$(\forall
a \in A)f(a) = g(a)$.
\end{proof}\noindent
ఇది సుక్రమాలు దృఢమైనవనే
భావాన్ని కొంత ఇస్తుంది.
దానిని మరింత వివరించడానికి
కొంత సంకేతలిపిని
పరిచయం చేయాలి.

\begin{defn}
$\tuple{A, <}$ ఒక సుక్రమం,
$a \in A$ అయితే,
$A_a = \Setabs{x \in A}{x
< a}$ అనుకుందాం.
$A_a$ను $A$ యొక్క
\emph{తనకన్నా చిన్న ఆరంభ ఖండం}
అంటాం ($A$ స్వయంగా
కూడా $A$ యొక్క చిన్నది కాని
ఆరంభ ఖండం కావచ్చని
అనుమతిస్తాం). $<_a$
అంటే $<$ను ఆరంభ
ఖండానికి పరిమితం
చేసిన సంబంధం; అంటే
$\funrestrictionto{\mathord{<}}{A_a^2}$.
\end{defn}
\noindent
ఈ సంకేతలిపితో, ఏ
సుక్రమమూ తనకన్నా చిన్న
ఆరంభ ఖండాల్లో దేనితోనూ
సమరూపం కాదని
ప్రకటించి నిరూపించవచ్చు.

\begin{lem}\ollabel{wellordnotinitial}
$\tuple{A, <}$ ఒక సుక్రమం,
$a \in A$ అయితే,
$\ordneq{\tuple{A, <}}{\tuple{A_a, <_a}}$.
\end{lem}

\begin{proof}
వైరుధ్యాన్ని చూపడానికి,
$f \colon A \to A_a$
సమరూపత అనుకుందాం.
$f$ ద్విగుణ ప్రమేయం,
$A_a \subsetneq A$ కాబట్టి,
\olref[wo]{wo:strictorder}ను
ఉపయోగించి $b \neq f(b)$
అయ్యే $A$లోని $<$-అత్యల్ప
\tetoken{మూలకం}{!!{element}}
$b \in A$ను తీసుకుందాం.
$(\forall x \in A)(x<b \liff x < f(b))$
అని చూపుతాం; అప్పుడు
\olref[sfr][rel][ord]{prop:extensionality-strictlinearorders}
ప్రకారం $b =
f(b)$ వచ్చి వైరుధ్యం
పూర్తవుతుంది.

$x < b$ అనుకుందాం.
$b$ ఎంపిక వల్ల
$x = f(x)$.
$f$ సమరూపత కాబట్టి
$f(x) <
f(b)$. అందువల్ల
$x < f(b)$.

$x < f(b)$ అనుకుందాం.
$f$ సమరూపత కాబట్టి
$f^{-1}(x) < b$;
$b$ ఎంపిక వల్ల
$f^{-1}(x) = x$.
అందువల్ల $x < b$.
\end{proof}

తదుపరి ఫలితం, స్థూలంగా,
సమరూపతలోని ``ఆరంభ
ఖండం'' కూడా సమరూపతే
అని చూపుతుంది:

\begin{lem}\ollabel{wellordinitialsegment}
$\tuple{A, <}$,
$\tuple{B, \lessdot}$
సుక్రమాలు అనుకుందాం.
$f
\colon A \to B$ సమరూపత,
$a \in A$ అయితే,
$\funrestrictionto{f}{A_{a}} : A_a \to B_{f(a)}$
సమరూపత అవుతుంది.
\end{lem}

\begin{proof}
$f$ సమరూపత కాబట్టి:
	%$f$ సమరూపత కాబట్టి, $b < a$ అప్పుడే మరియు అప్పుడే $f(b) \lessdot f(a)$; అందువల్ల
\begin{align*}
	\funimage{f}{A_a} &= \funimage{f}{\Setabs{x \in A}{x < a}}\\
	&= \funimage{f}{\Setabs{f^{-1}(y) \in A}{f^{-1}(y) < a}} \\
	&= \Setabs{y \in B}{y \lessdot f(a)} \\
	&=B_{f(a)} 
\end{align*}
ఇంకా $f$ క్రమాన్ని
సంరక్షిస్తుంది కాబట్టి
$\funrestrictionto{f}{A_a}$
కూడా సంరక్షిస్తుంది.
\end{proof}

తదుపరి రెండు ఫలితాలు
సుక్రమాలను ఎప్పుడూ
పోల్చవచ్చని స్థాపిస్తాయి:

\begin{lem}\ollabel{lemordsegments}
$\tuple{A, <}$,
$\tuple{B, \lessdot}$
సుక్రమాలు అనుకుందాం.
$\ordeq{\tuple{A_{a_1}, <_{a_1}}}{\tuple{B_{b_1}, \lessdot_{b_1}}}$
మరియు $\ordeq{\tuple{A_{{a_2}}, <_{a_2}}}{\tuple{B_{{b_2}},
\lessdot_{b_2}}}$ అయితే,
${a_1}  < {a_2} \text{ అప్పుడే మరియు అప్పుడే }{b_1} \lessdot
{b_2}$.
\end{lem}

\begin{proof}
\emph{ఎడమ నుంచి కుడికి}
నిరూపిస్తాం; మరో దిశ
కూడా ఇలాగే ఉంటుంది.
$\ordeq{\tuple{A_{a_1}, <_{a_1}}}{\tuple{B_{b_1},
\lessdot_{b_1}}}$, అలాగే
$\ordeq{\tuple{A_{{a_2}},
<_{a_2}}}{\tuple{B_{{b_2}}, \lessdot_{b_2}}}$
అనుకుందాం; మన సమరూపత
$f \colon
A_{{a_2}} \to B_{{b_2}}$.
${a_1} < {a_2}$ అనుకుందాం;
అప్పుడు
\olref{wellordinitialsegment}
ప్రకారం
$\ordeq{\tuple{A_{a_1}, <_{a_1}}}{\tuple{B_{f({a_1})},
\lessdot_{f({a_1})}}}$.
కాబట్టి
$\ordeq{\tuple{B_{b_1}, \lessdot_{b_1}}}{\tuple{B_{f({a_1})},
\lessdot_{f({a_1})}}}$;
అందువల్ల
\olref{wellordnotinitial}
ప్రకారం
${b_1} = f({a_1})$.
ఇప్పుడు $f$ యొక్క
విలువల పరిధి $B_{{b_2}}$
కాబట్టి ${b_1} \lessdot {b_2}$.
\sourcecorrection{OLTESTORDISO-001}{మూలంలో ఇక్కడ ``నిర్వచన సమితి'' అని ఉంది; సమరూపత లక్ష్య ఖండం కారణంగా ``విలువల పరిధి'' అని సరిచేశాం.}
\end{proof}

\begin{thm}\ollabel{thm:woalwayscomparable}
ఏ రెండు సుక్రమాలు ఇచ్చినా,
ఒకటి మరొకదాని
ఆరంభ ఖండంతో
(అది చిన్న ఖండమే
కావలసిన అవసరం లేదు)
సమరూపమవుతుంది.
\end{thm}

\begin{proof}
$\tuple{A, <}$,
$\tuple{B, \lessdot}$
సుక్రమాలు అనుకుందాం.
వేరుచేయడం ఉపయోగించి,
\[
	f = \Setabs{\tuple{a, b} \in A \times B}{
		\ordeq{\tuple{A_a, <_a}}{\tuple{B_b, \lessdot_b}}}.
\]
\olref{lemordsegments}
ప్రకారం ఇందులోని
అన్ని $\tuple{a_1, b_1}, \tuple{a_2, b_2} \in f$
కోసం $a_1 < a_2$
అప్పుడే మరియు అప్పుడే
$b_1 \lessdot b_2$.
కాబట్టి
$f \colon \dom{f} \to
\ran{f}$ ఒక సమరూపత.

$a_2 \in \dom{f}$,
$a_1 < a_2$ అయితే,
\olref{wellordinitialsegment}
ప్రకారం $a_1 \in \dom{f}$;
అందువల్ల $\dom{f}$,
$A$ యొక్క ఆరంభ
ఖండం. అదే విధంగా
$\ran{f}$, $B$
యొక్క ఆరంభ ఖండం.
వైరుధ్యాన్ని చూపడానికి
రెండూ తమ మూల
సమితులకన్నా \emph{చిన్న}
ఆరంభ ఖండాలే అనుకుందాం.
అప్పుడు $A \setminus \dom{f}$లోని
$<$-అత్యల్ప
\tetoken{మూలకం}{!!{element}}
$a$ను తీసుకుందాం;
అప్పుడు $\dom{f} =
A_a$. అలాగే
$B \setminus
\ran{f}$లోని
$\lessdot$-అత్యల్ప
\tetoken{మూలకం}{!!{element}}
$b$ను తీసుకుందాం;
అప్పుడు $\ran{f} = B_b$.
కాబట్టి
$f \colon A_a \to B_b$
సమరూపత; అందువల్ల
$\tuple{a, b} \in f$;
ఇది వైరుధ్యం.
\end{proof}

\end{document}
